% =============================================================================
% Low Power Edge Inference for Green Computing: 
% An Evaluation of Post-Training Quantization Techniques
% =============================================================================
% Conference Paper - Comprehensive Research Study
% Authors: Harshpreet Singh, Sameer Kumar
% Affiliation: Chitkara University, India
% =============================================================================

\documentclass[11pt,a4paper]{article}

% -----------------------------------------------------------------------------
% PAGE LAYOUT
% -----------------------------------------------------------------------------
\usepackage[margin=1in]{geometry}
\usepackage{setspace}
\doublespacing

% -----------------------------------------------------------------------------
% FONTS
% -----------------------------------------------------------------------------
\usepackage{times}
\usepackage[T1]{fontenc}

% -----------------------------------------------------------------------------
% MATH PACKAGES
% -----------------------------------------------------------------------------
\usepackage{amsmath}
\usepackage{amssymb}
\usepackage{amsthm}

% -----------------------------------------------------------------------------
% GRAPHICS AND TABLES
% -----------------------------------------------------------------------------
\usepackage{graphicx}
\usepackage{booktabs}
\usepackage{multirow}
\usepackage{array}
\usepackage{colortbl}
\usepackage{xcolor}
\usepackage{adjustbox}
\usepackage{tabularx}

% -----------------------------------------------------------------------------
% ALGORITHMS
% -----------------------------------------------------------------------------
\usepackage{algorithm}
\usepackage{algorithmic}

% -----------------------------------------------------------------------------
% FIGURES AND CAPTIONS
% -----------------------------------------------------------------------------
\usepackage{subcaption}
\usepackage{float}

% -----------------------------------------------------------------------------
% HYPERLINKS AND REFERENCES
% -----------------------------------------------------------------------------
\usepackage{hyperref}
\usepackage{url}
\hypersetup{
    colorlinks=true,
    linkcolor=blue,
    citecolor=blue,
    urlcolor=blue
}

% -----------------------------------------------------------------------------
% BIBLIOGRAPHY
% -----------------------------------------------------------------------------
\usepackage{natbib}
\bibliographystyle{plainnat}

% -----------------------------------------------------------------------------
% LISTS AND ENUMERATIONS
% -----------------------------------------------------------------------------
\usepackage{enumitem}
\setlist[itemize]{leftmargin=*,nosep}
\setlist[enumerate]{leftmargin=*,nosep}

% -----------------------------------------------------------------------------
% SYMBOLS AND SPECIAL CHARACTERS
% -----------------------------------------------------------------------------
\usepackage{pifont}
\usepackage{textcomp}

% -----------------------------------------------------------------------------
% PLOTS
% -----------------------------------------------------------------------------
\usepackage{pgfplots}
\pgfplotsset{compat=1.17}

% -----------------------------------------------------------------------------
% CUSTOM COMMANDS
% -----------------------------------------------------------------------------
\newcommand{\cmark}{\ding{51}}
\newcommand{\xmark}{\ding{55}}
\newcommand{\methodname}{PTQ-Edge}

% Math operators
\DeclareMathOperator*{\argmin}{arg\,min}
\DeclareMathOperator*{\argmax}{arg\,max}

% -----------------------------------------------------------------------------
% SECTION FORMATTING
% -----------------------------------------------------------------------------
\usepackage{titlesec}
\titleformat{\section}{\normalfont\Large\bfseries}{\thesection}{1em}{}
\titleformat{\subsection}{\normalfont\large\bfseries}{\thesubsection}{1em}{}
\titleformat{\subsubsection}{\normalfont\normalsize\bfseries}{\thesubsubsection}{1em}{}

% -----------------------------------------------------------------------------
% DOCUMENT METADATA
% -----------------------------------------------------------------------------
\title{\bfseries\Large Low Power Edge Inference for Green Computing: \\ 
An Evaluation of Post-Training Quantization on Apple Silicon}

\author{
  \textbf{Harshpreet Singh}$^{1}$ \and \textbf{Sameer Kumar}$^{1}$ \\[6pt]
  $^{1}$Chitkara University, India \\[3pt]
  \texttt{\{harshpreet393.be22, sameer776.be22\}@chitkara.edu.in}
}

\date{}

\begin{document}

\maketitle

% =============================================================================
% ABSTRACT
% =============================================================================
\begin{abstract}
\noindent We present a comprehensive evaluation of post-training quantization (PTQ) techniques for deploying large language models on resource-constrained edge devices, addressing the critical need for sustainable AI inference. With global data center energy consumption projected to reach 945 TWh annually by 2030—driven primarily by AI workloads—edge computing offers a promising pathway to reduce both operational costs and environmental impact. We systematically evaluate inference frameworks across two model families: Qwen3.5 series (0.8B, 2B, and 4B parameters) and Microsoft Phi-3 Mini (3.82B parameters). Our experiments on Apple Silicon M4 Pro with 24GB unified memory reveal that Apple's MLX framework with INT4 quantization achieves 269.88, 180.48, 80.69, and 96.53 tokens/second for Qwen3.5-0.8B, Qwen3.5-2B, Qwen3.5-4B, and Phi-3 Mini respectively, with memory footprints ranging from 0.48GB to 2.47GB. Comparing frameworks, MLX provides up to 9.9$\times$ speedup over Ollama for Qwen3.5-0.8B (269.88 vs 27.24 tokens/second). Our findings provide actionable guidance for practitioners: Qwen3.5-0.8B with MLX INT4 quantization offers exceptional throughput for real-time applications, while all tested models fit comfortably within 8GB memory, enabling deployment on consumer devices.
\end{abstract}

% =============================================================================
% 1. INTRODUCTION
% =============================================================================
\section{Introduction}

The proliferation of large language models (LLMs) has transformed natural language processing, enabling unprecedented capabilities in text generation, reasoning, and knowledge retrieval \cite{brown2020language, openai2023gpt4}. However, this progress comes at a substantial environmental cost. According to the International Energy Agency, data center electricity consumption is projected to reach 945 terawatt-hours (TWh) annually by 2030, representing approximately 3\% of global electricity consumption—with AI workloads driving the majority of this growth \cite{iea2025energyai}. This exponential increase in energy demand underscores an urgent need for sustainable AI deployment strategies.

\subsection{The Sustainability Imperative}

The environmental impact of LLMs extends across their entire lifecycle. Training a single large language model can emit as much carbon dioxide as five cars over their lifetimes \cite{strubell2019energy}, while inference—performed billions of times daily—accounts for the majority of operational energy consumption in production deployments \cite{patterson2021carbon}. As models grow larger to achieve better performance, the energy required for both training and inference scales accordingly, creating a fundamental tension between capability and sustainability.

Edge computing presents a compelling alternative to cloud-based inference, offering reduced latency, enhanced privacy, and potential energy savings through local computation \cite{shi2016edge}. However, deploying LLMs on edge devices introduces significant challenges: limited memory capacity, constrained computational resources, and strict power budgets. A 4-billion parameter model in FP16 precision requires approximately 8GB of memory for weights alone—exceeding the capabilities of many edge devices and demanding compression techniques that preserve model capability while reducing resource requirements.

\subsection{The Quantization Challenge}

Post-training quantization (PTQ) has emerged as a leading technique for model compression, reducing the numerical precision of model weights from 16-bit or 32-bit floating-point to lower-bit representations such as 8-bit or 4-bit integers \cite{han2015deep}. Unlike quantization-aware training (QAT), which requires retraining the model, PTQ operates on pre-trained weights without additional training—making it particularly attractive for deploying existing models on resource-constrained hardware.

However, the landscape of PTQ methods is fragmented, with different techniques exhibiting varying characteristics across different model architectures. Key questions remain unanswered:

\begin{itemize}[leftmargin=*,nosep]
    \item Which PTQ method provides the best memory efficiency for models of different sizes and architectures?
    \item How do different quantization strategies affect energy consumption during inference?
    \item What is the optimal bit-width for balancing memory efficiency and inference speed?
    \item How does model architecture influence quantization sensitivity?
\end{itemize}

\subsection{Contributions}

This paper addresses these questions through a systematic evaluation of PTQ techniques for edge inference. Our contributions are:

\begin{enumerate}[leftmargin=*,nosep]
    \item \textbf{Framework Comparison}: We evaluate Ollama and MLX frameworks across Qwen3.5 series (0.8B, 2B, 4B parameters) and Phi-3 Mini (3.82B parameters), finding that MLX provides up to 9.9$\times$ speedup over Ollama for Qwen3.5-0.8B (269.88 vs 27.24 tokens/second).
    
    \item \textbf{Real Measurements}: We measure actual inference latency (3.71, 5.54, 12.39, and 10.36 ms/token for Qwen3.5-0.8B, 2B, 4B, and Phi-3 Mini via MLX INT4), energy consumption (0.000053 to 0.000361 kWh), and memory footprint (0.48GB to 2.47GB with MLX INT4).
    
    \item \textbf{Practical Guidelines}: We provide actionable recommendations: Qwen3.5-0.8B with MLX INT4 quantization achieves exceptional throughput (269.88 tokens/second) for real-time applications, while all tested models fit within 8GB memory for consumer device deployment.
\end{enumerate}

Our findings have immediate practical implications for sustainable AI deployment, enabling organizations to reduce their inference energy footprint while maintaining model capability.

% =============================================================================
% 2. BACKGROUND
% =============================================================================
\section{Background and Motivation}

\subsection{The Energy Crisis in AI Computing}

The intersection of artificial intelligence and environmental sustainability has become a critical research frontier. The computational demands of modern LLMs have grown exponentially: while BERT-base (110M parameters, 2018) could be trained on a single GPU in hours, contemporary models like GPT-4 (estimated 1.8T parameters) require thousands of GPUs running for months \cite{openai2023gpt4}. This scaling has direct environmental consequences.

\subsubsection{Data Center Energy Projections}

The International Energy Agency's 2025 report on Energy and AI provides sobering statistics \cite{iea2025energyai}:

\begin{itemize}[leftmargin=*,nosep]
    \item Data centers consumed approximately 460 TWh globally in 2024, representing 1.5\% of global electricity consumption
    \item By 2030, this is projected to reach 945 TWh (3\% of global consumption), with sensitivity scenarios ranging from 700 TWh to 1,700 TWh
    \item AI workloads are the primary driver, with server power density increasing from 5-10 kW/rack to 20-40 kW/rack for AI-optimized infrastructure
    \item The United States and China account for nearly 80\% of global data center energy consumption growth
\end{itemize}

These projections assume continued efficiency improvements in hardware and cooling; without such improvements, energy consumption would be significantly higher.

\subsubsection{Inference vs. Training Energy}

While training receives significant attention due to its concentrated energy expenditure, inference dominates operational energy consumption over a model's lifetime. Patterson et al. \cite{patterson2021carbon} estimate that inference accounts for approximately 60\% of total ML compute usage in production systems, with this percentage increasing for widely-deployed models that serve billions of requests.

For a model serving 1 billion queries annually, the cumulative inference energy can exceed training energy by an order of magnitude. This makes inference efficiency a critical lever for reducing AI's environmental footprint.

\subsection{Edge Computing for Sustainable AI}

Edge computing—the deployment of computational resources at the network periphery, close to data sources and end users—offers several advantages for sustainable AI deployment \cite{shi2016edge}:

\begin{enumerate}[leftmargin=*,nosep]
    \item \textbf{Reduced Data Transfer}: Processing data locally eliminates the energy cost of transmitting data to and from cloud data centers, which can account for 30-40\% of total energy in cloud-based inference pipelines \cite{masanet2020recalibrating}.
    
    \item \textbf{Utilization of Existing Hardware}: Edge devices are already deployed and powered; utilizing their computational capacity for AI inference represents a more efficient use of embodied carbon compared to additional cloud infrastructure.
    
    \item \textbf{Reduced Cooling Overhead}: Edge devices typically operate in ambient environments without dedicated cooling infrastructure, whereas data centers require 30-40\% of their energy for cooling \cite{avgerinou2017towards}.
    
    \item \textbf{Privacy and Latency}: Edge inference eliminates privacy concerns associated with transmitting user data to cloud servers and reduces latency for time-sensitive applications.
\end{enumerate}

However, edge devices present significant constraints for LLM deployment:

\begin{itemize}[leftmargin=*,nosep]
    \item \textbf{Memory Constraints}: Consumer edge devices typically have 8-16GB of unified memory, while LLMs require memory proportional to parameter count (2 bytes per parameter in FP16).
    
    \item \textbf{Compute Limitations}: Edge processors have lower peak FLOPs than data center GPUs, making efficient inference critical.
    
    \item \textbf{Power Budgets}: Battery-powered devices have strict power constraints that limit sustained computation.
\end{itemize}

These constraints motivate the need for effective model compression techniques.

\subsection{Model Compression Techniques}

Model compression encompasses a family of techniques for reducing the memory footprint and computational cost of neural networks while preserving their predictive capability \cite{cheng2017survey}. The primary approaches include:

\subsubsection{Quantization}

Quantization reduces the numerical precision of model parameters, mapping continuous floating-point values to discrete low-bit representations. For a weight tensor $\mathbf{W} \in \mathbb{R}^{n \times m}$, quantization to $b$ bits involves:

\begin{equation}
\mathbf{W}_q = \text{clamp}\left(\text{round}\left(\frac{\mathbf{W}}{s}\right), -2^{b-1}, 2^{b-1}-1\right)
\end{equation}

where $s$ is the scale factor determining the quantization step size. The dequantized approximation is:

\begin{equation}
\hat{\mathbf{W}} = s \cdot \mathbf{W}_q
\end{equation}

The quantization error $\|\mathbf{W} - \hat{\mathbf{W}}\|$ directly impacts model accuracy, and different PTQ methods employ various strategies to minimize this error.

\subsubsection{Pruning}

Pruning removes redundant parameters from the network, setting weights to zero based on importance criteria. Unstructured pruning operates on individual weights:

\begin{equation}
\mathbf{W}_{\text{pruned}} = \mathbf{W} \odot \mathbf{M}, \quad \mathbf{M}_{ij} = \begin{cases} 1 & \text{if } |W_{ij}| > \tau \\ 0 & \text{otherwise} \end{cases}
\end{equation}

where $\tau$ is a threshold and $\odot$ denotes element-wise multiplication. The sparsity level is controlled by adjusting $\tau$.

\subsubsection{Knowledge Distillation}

Knowledge distillation trains a smaller ``student'' model to mimic a larger ``teacher'' model \cite{hinton2015distilling}. While effective, this requires training access to the teacher model and significant computational resources, making it less suitable for post-training compression scenarios.

\subsection{The Case for PTQ in Edge Deployment}

Post-training quantization offers several advantages for edge AI deployment:

\begin{enumerate}[leftmargin=*,nosep]
    \item \textbf{No Training Required}: PTQ operates directly on pre-trained weights, making it applicable to any existing model without access to training infrastructure or data.
    
    \item \textbf{Rapid Deployment}: Quantization can be performed in minutes to hours (depending on model size), compared to weeks for retraining-based approaches.
    
    \item \textbf{Hardware Compatibility}: Many edge processors include specialized hardware for low-precision arithmetic (e.g., Apple Neural Engine, Qualcomm Hexagon DSP), accelerating quantized inference.
    
    \item \textbf{Memory Efficiency}: 4-bit quantization reduces memory footprint by 4$\times$ compared to FP16, enabling deployment of larger models on memory-constrained devices.
\end{enumerate}

These advantages make PTQ the most practical approach for sustainable edge AI deployment, motivating our systematic evaluation of PTQ methods.

% =============================================================================
% 3. RELATED WORK
% =============================================================================
\section{Related Work}

\subsection{Post-Training Quantization Methods}

The development of PTQ methods has evolved rapidly, driven by the need to compress increasingly large language models. We categorize the major approaches:

\subsubsection{INT8 Uniform Quantization}

Uniform quantization maps floating-point values to uniformly spaced integer levels. For LLMs, LLM.int8() \cite{dettmers2022llm} introduced mixed-precision decomposition, identifying outlier feature dimensions that require higher precision:

\begin{equation}
\mathbf{Y} = \mathbf{X}_{\text{normal}} \mathbf{W}_{\text{int8}} + \mathbf{X}_{\text{outlier}} \mathbf{W}_{\text{fp16}}
\end{equation}

where $\mathbf{X}_{\text{outlier}}$ contains dimensions with values exceeding a threshold. This approach preserves accuracy for most models while providing memory savings, though with increased computational overhead for the decomposition.

Simple INT8 weight-only quantization provides consistent 2$\times$ memory reduction with minimal accuracy loss for most models, making it a reliable baseline approach.

\subsubsection{GPTQ: Optimal Brain Quantization}

GPTQ \cite{frantar2022gptq} adapts Optimal Brain Surgeon (OBS) \cite{lecun1990optimal} for LLM quantization. The key insight is to quantize weights sequentially while updating remaining weights to compensate for quantization error:

\begin{equation}
\delta_{\mathbf{W}} = -\frac{w_q - w}{[\mathbf{H}^{-1}]_{qq}} \cdot \mathbf{H}^{-1}_{:,q}
\end{equation}

where $\mathbf{H}$ is the Hessian matrix of the layer's output, $w$ is the original weight, $w_q$ is the quantized weight, and $q$ indexes the quantized weight position. The update $\delta_{\mathbf{W}}$ adjusts unquantized weights to minimize output error.

GPTQ achieves impressive results, enabling 4-bit quantization of 175B parameter models with minimal perplexity degradation. However, the algorithm requires computing the inverse Hessian, which has $O(d^3)$ complexity for a layer with $d$ inputs—addressed through Cholesky decomposition and lazy batch updates.

\subsubsection{AWQ: Activation-Aware Weight Quantization}

AWQ \cite{lin2023awq} observes that not all weights are equally important: protecting just 1\% of salient weights can significantly reduce quantization error. The key innovation is identifying salient weights based on activation magnitudes:

\begin{equation}
s^* = \argmin_{s} \mathcal{L}(\mathbf{W}, s), \quad \mathcal{L} = \sum_{i,j} \left| \mathbf{X}_{i,j} \cdot (w_{i,j} - \hat{w}_{i,j}) \right|
\end{equation}

where $s$ is a per-channel scaling factor, $\mathbf{X}$ represents activations from calibration data, and the scaling is applied as:

\begin{equation}
\hat{\mathbf{W}} = \text{quant}(\mathbf{W} \cdot \text{diag}(s)) \cdot \text{diag}(s)^{-1}
\end{equation}

This activation-aware scaling preserves salient weights during quantization, achieving near-lossless 4-bit compression for many models. AWQ received the Best Paper Award at MLSys 2024.

\subsubsection{INT4 Weight-Only Quantization}

Naive INT4 weight-only quantization applies uniform quantization to 4-bit precision without sophisticated error compensation. While this provides the maximum memory savings among common PTQ methods, it typically incurs significant accuracy degradation, especially for larger models or those with complex weight distributions.

However, INT4 quantization serves as an important baseline and is increasingly viable for models specifically designed for low-precision deployment.

\subsection{Quantization for Edge Deployment}

Several studies have explored LLM quantization for edge scenarios:

\textbf{llama.cpp} \cite{llamacpp2023} introduced GGUF format with various quantization schemes (Q4\_0, Q4\_K\_M, etc.) optimized for CPU inference. Their block quantization approach balances compression and accuracy.

\textbf{MLC-LLM} \cite{mlcllm2023} provides a compilation-based approach for deploying quantized LLMs across diverse hardware platforms, including mobile devices and Apple Silicon.

\textbf{Huang et al.} \cite{huang2025llms} evaluated LLM inference on edge devices using Ollama, measuring energy efficiency across 28 quantized models on Raspberry Pi 4.

\textbf{Benazir and Lin} \cite{benazir2025profiling} profiled LLM inference on Apple Silicon from a quantization perspective, identifying optimal configurations for the M-series architecture.

Our work extends these studies by providing a systematic comparison of inference frameworks (Ollama and MLX) specifically for recent small language models (Qwen3.5 series and Phi-3 Mini) with energy measurements on Apple Silicon.

\subsection{Small Language Models}

Small language models (SLMs) have emerged as efficient alternatives to large models for edge deployment:

\textbf{Phi Series} \cite{javaheripi2023phi2, abdin2024phi3}: Microsoft's Phi models demonstrate that carefully curated training data can enable small models to match the performance of much larger models. Phi-3 Mini (3.82B parameters) achieves strong performance through textbook-quality training data while maintaining a compact footprint suitable for edge deployment.

\textbf{Qwen Series} \cite{qwen2025technical}: Alibaba's Qwen models include dense and Mixture-of-Expert architectures from 0.5B to 72B parameters. The Qwen3.5 series introduces hybrid reasoning with thinking and non-thinking modes, achieving competitive performance through architectural innovations including GQA attention and dual-mode inference.

These models are particularly relevant for edge deployment due to their optimized size-performance trade-off, though their quantization characteristics remain understudied.

\subsection{Energy Measurement in ML}

Measuring energy consumption for ML inference requires careful methodology:

\textbf{CodeCarbon} \cite{codecarbon2023}: A Python package that tracks emissions by monitoring CPU, GPU, and RAM power consumption, integrating with regional carbon intensity data. It supports both online and offline modes.

\textbf{RAPL (Running Average Power Limit)}: Intel's interface for measuring processor power consumption at fine granularity, widely used in ML energy studies.

\textbf{Hardware Power Meters}: External power meters provide ground-truth measurements but require physical instrumentation.

For Apple Silicon, energy measurement presents unique challenges due to the unified memory architecture and limited software interfaces. We address this through Apple's powerlog framework and CodeCarbon's macOS support.

% =============================================================================
% 4. METHODOLOGY
% =============================================================================
\section{Methodology}

\subsection{Experimental Design}

Our experimental design follows a systematic approach to evaluate PTQ methods on Apple Silicon:

\begin{enumerate}[leftmargin=*,nosep]
    \item \textbf{Framework Comparison}: Evaluate Ollama and MLX frameworks for inference optimization across all model families, comparing quantization strategies and inference speeds.
    \item \textbf{Model Selection}: Test on Qwen3.5 series (0.8B, 2B, 4B) and Phi-3 Mini, representing diverse size-efficiency trade-offs.
    \item \textbf{Comprehensive Metrics}: Measure latency, memory footprint, and energy consumption.
\end{enumerate}

\subsection{Models}

\subsubsection{Qwen3.5 Series}

The Qwen3.5 series \cite{qwen2025technical} from Alibaba represents the latest generation of efficient language models with hybrid reasoning capabilities. We evaluate three variants:

\begin{itemize}[leftmargin=*,nosep]
    \item \textbf{Qwen3.5-0.8B}: 873M parameters, optimized for extreme efficiency
    \item \textbf{Qwen3.5-2B}: 2.27B parameters, balanced size-performance trade-off
    \item \textbf{Qwen3.5-4B}: 4.66B parameters, highest capability in the series
\end{itemize}

Key architectural features shared across the series:
\begin{itemize}[leftmargin=*,nosep]
    \item \textbf{Architecture}: Dense transformer with GQA (Grouped Query Attention)
    \item \textbf{Position Encoding}: RoPE (Rotary Position Embedding)
    \item \textbf{Activation}: SwiGLU
    \item \textbf{Hybrid Reasoning}: Dual-mode inference (thinking and non-thinking modes)
    \item \textbf{Context Length}: Up to 32,768 tokens
    \item \textbf{Languages}: Supports 119 languages
\end{itemize}

The thinking mode enables extended reasoning for complex tasks, while non-thinking mode provides fast responses for simpler queries. We evaluate Qwen3.5 models via both Ollama and MLX frameworks.

\subsubsection{Phi-3 Mini}

Phi-3 Mini \cite{abdin2024phi3} is a 3.82-billion parameter model from Microsoft's Phi-3 family, designed for efficient edge deployment:

\begin{itemize}[leftmargin=*,nosep]
    \item \textbf{Architecture}: Dense transformer with optimized attention
    \item \textbf{Hidden Dimensions}: 3072
    \item \textbf{Attention Heads}: 32
    \item \textbf{Layers}: 32
    \item \textbf{Context Length}: 128,000 tokens (with RoPE)
    \item \textbf{Training}: Trained on 3.3T tokens with high-quality educational data
\end{itemize}

Phi-3 Mini is available through both Ollama (Q4\_K\_M quantization) and MLX (INT4 quantization). In our experiments, MLX INT4 quantization achieved 2.29GB memory footprint with 96.53 tokens/second inference speed.


\subsection{Hardware Platform}

All experiments were conducted on:

\begin{itemize}[leftmargin=*,nosep]
    \item \textbf{Hardware}: Apple Silicon M4 Pro (14-core CPU, 20-core GPU)
    \item \textbf{Memory}: 24GB unified memory
    \item \textbf{Software}: macOS, PyTorch 2.7.1 with MPS backend
\end{itemize}


\subsection{Evaluation Metrics}


\subsubsection{Latency}

We measure inference latency for generating 128 tokens with 256-token input context:

\begin{equation}
\text{Latency} = \frac{\text{Total Generation Time}}{\text{Number of Tokens}}
\end{equation}

We report mean and standard deviation over 10 runs after warmup.

\subsubsection{Memory Footprint}

We measure peak memory consumption during inference, including:

\begin{itemize}[leftmargin=*,nosep]
    \item Model weights (quantized or full-precision)
    \item KV cache for attention
    \item Activation buffers
    \item Runtime overhead
\end{itemize}

\subsubsection{Energy Consumption}

We measure energy using CodeCarbon \cite{codecarbon2023} integrated with Apple's power management framework:

\begin{equation}
E_{\text{inference}} = \int_{t_0}^{t_1} P(t) \, dt
\end{equation}

where $P(t)$ is instantaneous power consumption. We report energy in Joules per inference and compute energy efficiency as tokens per Joule.

\subsection{Experimental Setup}

\subsubsection{Hardware}

All experiments are conducted on:
\begin{itemize}[leftmargin=*,nosep]
    \item \textbf{Platform}: Apple MacBook Pro (M4 Pro)
    \item \textbf{Processor}: Apple M4 Pro (12-core CPU, 16-core GPU)
    \item \textbf{Memory}: 24GB unified memory (LPDDR5X)
    \item \textbf{Storage}: 512GB SSD
    \item \textbf{OS}: macOS Sequoia 15.3
\end{itemize}

The Apple Silicon architecture provides unified memory accessible by both CPU and GPU, eliminating data transfer overhead and enabling efficient inference for quantized models.

\subsubsection{Software}

\begin{itemize}[leftmargin=*,nosep]
    \item \textbf{Framework}: PyTorch 2.5 with Metal Performance Shaders (MPS)
    \item \textbf{Quantization}: MLX framework
    \item \textbf{Inference}: llama.cpp (GGUF format), MLX
    \item \textbf{Energy Measurement}: CodeCarbon 3.2.3
\end{itemize}


\subsubsection{Hyperparameters}

\begin{itemize}[leftmargin=*,nosep]
    \item Calibration samples: 512
    \item Calibration dataset: The Pile (random subset)
    \item Generation tokens: 128
    \item Input context: 256 tokens
    \item Temperature: 0.0 (greedy decoding)
    \item Random seeds: [42, 123, 456] for reproducibility
\end{itemize}

% =============================================================================
% 5. EXPERIMENTS AND RESULTS
% =============================================================================
\section{Experiments and Results}


\subsection{Memory Footprint}

Table \ref{tab:memory} shows measured memory requirements from experiments on Apple Silicon M4 Pro.

\begin{table}[t]
\centering
\caption{Model specifications and memory requirements on Apple Silicon M4 Pro with 24GB unified memory. Qwen3.5 models evaluated via Ollama with Q4\_K\_M quantization; Phi-3 Mini evaluated via MLX with INT4 quantization.}
\label{tab:memory}
\begin{tabular}{l|ccc}
\toprule
\textbf{Model} & \textbf{Parameters} & \textbf{Size (GB)} & \textbf{Framework} \\
\midrule
Qwen3.5-0.8B & 873M & 0.52 & Ollama Q4\_K\_M \\
Qwen3.5-2B & 2.27B & 1.32 & Ollama Q4\_K\_M \\
Qwen3.5-4B & 4.66B & 2.65 & Ollama Q4\_K\_M \\
Phi-3 Mini & 3.82B & 2.28 & MLX INT4 \\
\bottomrule
\end{tabular}
\end{table}

The model sizes are measured directly from loaded weights in their respective quantized formats. With MLX 0.31.0, Qwen3.5 models are now fully supported with INT4 quantization.

\subsection{Latency Analysis}

Table \ref{tab:latency} presents inference latency measurements from real experiments.

\begin{table}[t]
\centering
\caption{Inference latency measurements on Apple Silicon M4 Pro. Qwen3.5 models via Ollama (Q4\_K\_M quantization, non-thinking mode); Phi-3 Mini via MLX (INT4 quantization).}
\label{tab:latency}
\begin{tabular}{l|cccc}
\toprule
\textbf{Metric} & \textbf{Qwen3.5-0.8B} & \textbf{Qwen3.5-2B} & \textbf{Qwen3.5-4B} & \textbf{Phi-3 Mini} \\
\midrule
ms/token & 36.71 & 55.33 & 188.78 & 10.36 \\
tokens/sec & 27.24 & 18.07 & 5.30 & 96.53 \\
Framework & Ollama & Ollama & Ollama & MLX \\
\bottomrule
\end{tabular}
\end{table}
The results show clear scaling behavior: smaller models achieve significantly higher throughput. Qwen3.5-0.8B achieves 27.24 tokens/second via Ollama, while Phi-3 Mini with MLX INT4 quantization achieves an impressive 96.53 tokens/second due to hardware-optimized inference.

\subsection{MLX Quantization Results}

Apple's MLX framework provides hardware-optimized quantization for Apple Silicon. With MLX 0.31.0, we successfully benchmarked Qwen3.5 models with INT4 quantization. Table \ref{tab:mlx_memory} presents MLX INT4 quantization results for all models.

\begin{table}[t]
\centering
\caption{MLX INT4 quantization results on Apple Silicon M4 Pro. All models evaluated with INT4 quantization. Qwen3.5-0.8B achieves the highest throughput at 269.88 tokens/second with only 0.48GB memory.}
\label{tab:mlx_memory}
\begin{tabular}{l|cccc}
\toprule
\textbf{Model} & \textbf{Quantization} & \textbf{Memory (GB)} & \textbf{ms/token} & \textbf{tokens/sec} \\
\midrule
Qwen3.5-0.8B & INT4 & 0.48 & 3.71 & 269.88 \\
Qwen3.5-2B & INT4 & 1.12 & 5.54 & 180.48 \\
Qwen3.5-4B & INT4 & 2.47 & 12.39 & 80.69 \\
Phi-3 Mini & INT4 & 2.29 & 10.36 & 96.53 \\
\bottomrule
\end{tabular}
\end{table}

Key findings from MLX quantization:

\begin{enumerate}[leftmargin=*,nosep]
    \item \textbf{Dramatic Speedup}: MLX INT4 quantization provides a 9.9$\times$ speedup for Qwen3.5-0.8B compared to Ollama (269.88 vs 27.24 tokens/second), demonstrating the importance of hardware-optimized inference frameworks.
    
    \item \textbf{Memory Efficiency}: All models fit comfortably under 2.5GB memory with INT4 quantization, with Qwen3.5-0.8B requiring only 0.48GB, enabling deployment on devices with limited RAM.
    
    \item \textbf{Framework Selection}: For Apple Silicon, MLX significantly outperforms Ollama across all model sizes, with speedups ranging from 7.4$\times$ (Phi-3 Mini: 96.53 vs 13.04 tok/s) to 15.2$\times$ (Qwen3.5-4B: 80.69 vs 5.30 tok/s).
    
    \item \textbf{Practical Implications}: MLX enables real-time inference for all tested models on devices with 8GB memory, with Qwen3.5-0.8B achieving near-instantaneous response times at 269.88 tokens/second.
\end{enumerate}
\subsection{Ollama Inference Results}

We benchmarked Qwen3.5 series models using Ollama with Q4\_K\_M quantization (default for Ollama models). Table~\ref{tab:ollama_results} presents inference results on Apple Silicon M4 Pro in non-thinking mode.

\begin{table}[t]
\centering
\caption{Ollama inference results on Apple Silicon M4 Pro with Q4\_K\_M quantization (non-thinking mode). Energy measured using CodeCarbon.}
\label{tab:ollama_results}
\begin{tabular}{l|cccc}
\toprule
\textbf{Model} & \textbf{ms/token} & \textbf{tok/s} & \textbf{Energy (kWh)} & \textbf{Parameters} \\
\midrule
Qwen3.5-0.8B & 36.71 & 27.24 & 0.000053 & 873M \\
Qwen3.5-2B & 55.33 & 18.07 & 0.000141 & 2.27B \\
Qwen3.5-4B & 188.78 & 5.30 & 0.000361 & 4.66B \\
Phi-3 Mini & 76.68 & 13.04 & 0.000100 & 3.82B \\
\bottomrule
\end{tabular}
\end{table}
The results demonstrate clear scaling behavior: Qwen3.5-0.8B achieves the highest throughput (27.24 tokens/second) with the lowest energy consumption (0.000053 kWh), making it ideal for energy-constrained edge deployment. As model size increases, latency increases and throughput decreases proportionally. Phi-3 Mini provides a middle ground with 13.04 tokens/second, suitable for applications requiring more capability than the smallest Qwen3.5 variant.


\subsection{Energy Consumption}

Table \ref{tab:energy} shows energy consumption measurements tracked by CodeCarbon during inference experiments.

\begin{table}[t]
\centering
\caption{Energy consumption per inference (generating 100 tokens) measured using CodeCarbon on Apple Silicon M4 Pro. Qwen3.5 models via Ollama; Phi-3 Mini via Ollama.
\label{tab:energy}
\begin{tabular}{l|ccc}
\toprule
\textbf{Model} & \textbf{Energy (kWh)} & \textbf{Energy (Joules)} & \textbf{Efficiency (tok/J)} \\
\midrule
Qwen3.5-0.8B & 0.000053 & 190.8 & 524 \\
Qwen3.5-2B & 0.000141 & 507.6 & 197 \\
Qwen3.5-4B & 0.000361 & 1299.6 & 77 \\
Phi-3 Mini & 0.000100 & 360.0 & 278 \\
\bottomrule
\end{tabular}
\end{table}
Energy efficiency (tokens per Joule) decreases with model size, with Qwen3.5-0.8B achieving 524 tokens/Joule compared to 77 tokens/Joule for Qwen3.5-4B. Phi-3 Mini achieves 278 tokens/Joule, providing a balance between capability and efficiency. These measurements highlight the importance of model selection for energy-constrained edge deployment.

\subsection{Experimental Limitations}

Our experiments revealed several important limitations for edge inference on Apple Silicon:

\begin{enumerate}[leftmargin=*,nosep]
    
    
    \item \textbf{Energy measurement}: CodeCarbon uses estimated power consumption for Apple Silicon. More accurate measurement would require external power meters.
    
    \item \textbf{Framework comparison}: MLX provides significant speedups (7.4$\times$ to 15.2$\times$) over Ollama across all tested models, demonstrating the importance of hardware-optimized inference on Apple Silicon.
\end{enumerate}
\subsection{Recommendations for Practitioners}

Based on our experimental findings:

\begin{enumerate}[leftmargin=*,nosep]
    \item \textbf{Model selection}: Qwen3.5-0.8B with MLX INT4 quantization provides exceptional throughput (269.88 tok/s) for real-time applications requiring the fastest response times.
    
    \item \textbf{Framework choice}: Use MLX for all models on Apple Silicon to achieve optimal performance, with speedups ranging from 7.4$\times$ (Phi-3 Mini) to 15.2$\times$ (Qwen3.5-4B) compared to Ollama.
    
    \item \textbf{Memory constraints}: All tested models with MLX INT4 quantization fit within 8GB memory (0.48GB to 2.47GB), enabling deployment on consumer devices.
\end{enumerate}


% =============================================================================
% 6. DISCUSSION

Our experiments on Apple Silicon M4 Pro reveal important insights for edge LLM deployment:



\subsubsection{Model-Specific Observations}

\begin{itemize}[leftmargin=*,nosep]
    \item \textbf{Qwen3.5 Series}: With MLX INT4 quantization, dramatic speedups are observed across all sizes. Qwen3.5-0.8B achieves 269.88 tokens/second (3.71 ms/token), representing a 9.9$\times$ speedup over Ollama. Qwen3.5-2B achieves 180.48 tokens/second, and Qwen3.5-4B achieves 80.69 tokens/second, making all variants suitable for real-time edge applications.
    
    \item \textbf{Phi-3 Mini}: Achieves 13.04 tokens/second via Ollama and 96.53 tokens/second via MLX INT4 quantization, demonstrating a 7.4$\times$ speedup and the importance of framework selection on Apple Silicon.
\end{itemize}
\subsubsection{Energy Measurement Challenges}

CodeCarbon's constant power model for Apple Silicon provides estimates rather than precise measurements. Key energy findings for inference:

\begin{itemize}[leftmargin=*,nosep]
    \item Qwen3.5-0.8B: 0.000053 kWh (most efficient at 524 tokens/Joule)
    \item Qwen3.5-2B: 0.000141 kWh (197 tokens/Joule)
    \item Qwen3.5-4B: 0.000361 kWh (77 tokens/Joule)
    \item Phi-3 Mini: 0.000100 kWh (278 tokens/Joule)
    \item Energy scales approximately linearly with parameter count
\end{itemize}
% =============================================================================
% 6. DISCUSSION
% =============================================================================
\section{Discussion}

\subsection{Implications for Green AI}

Our experiments on Apple Silicon M4 Pro provide valuable insights for edge LLM deployment:

\subsubsection{Hardware-Specific Considerations}

Apple Silicon's unified memory architecture offers advantages for LLM inference:

\begin{itemize}[leftmargin=*,nosep]
    \item No GPU-CPU data transfer bottleneck
    \item 24GB unified memory accommodates all tested models in quantized form
    \item Metal Performance Shaders enable efficient inference (96.53 tok/s for Phi-3 Mini INT4, up to 269.88 tok/s for Qwen3.5-0.8B INT4)
    \item Framework choice significantly impacts performance (MLX 7.4$\times$ faster than Ollama for Phi-3 Mini, up to 15.2$\times$ for Qwen3.5-4B)
\end{itemize}


\subsubsection{Energy Consumption Patterns}

Energy efficiency varies significantly with model size:

\begin{itemize}[leftmargin=*,nosep]
    \item Smallest model (Qwen3.5-0.8B) achieves 524 tokens/Joule
    \item Largest model (Qwen3.5-4B) achieves only 77 tokens/Joule
    \item Phi-3 Mini balances capability and efficiency at 278 tokens/Joule
    \item Energy scales approximately linearly with parameter count
\end{itemize}
\subsection{Limitations}



\subsubsection{Hardware Specificity}

Results are specific to Apple Silicon M4 Pro:

\begin{itemize}[leftmargin=*,nosep]
    \item NVIDIA GPUs may show different memory savings
    \item Mobile NPUs have different quantization support
    \item Results may not generalize to other Apple Silicon variants
\end{itemize}

\subsubsection{Model Selection}

We evaluate Qwen3.5 series and Phi-3 Mini; generalization to other architectures requires further study:

\begin{itemize}[leftmargin=*,nosep]
    \item Mixture-of-Expert models may have different quantization characteristics
    \item Qwen3.5 thinking mode (extended reasoning) not evaluated in this study
    \item Models with different training procedures may respond differently
    \item Very large models (>10B) were not evaluated
\end{itemize}


\subsection{Future Work}

Several directions warrant further investigation:

\begin{enumerate}[leftmargin=*,nosep]
    \item \textbf{Cross-Architecture Evaluation}: Extend experiments to NVIDIA GPUs, mobile NPUs, and embedded processors.
    
    \item \textbf{Dynamic Quantization}: Explore dynamic bit-width selection based on layer sensitivity.
    
    \item \textbf{Task-Specific Optimization}: Investigate whether different tasks benefit from different quantization strategies.
    
    \item \textbf{Combined Techniques}: Systematic study of quantization combined with other compression (distillation, structured pruning).
    
    \item \textbf{Carbon-Aware Deployment}: Integrate quantization decisions with real-time carbon intensity data for truly green inference.
\end{enumerate}

% =============================================================================
% 7. CONCLUSION
% =============================================================================
\section{Conclusion}

This paper presents an experimental evaluation of inference frameworks for deploying language models on Apple Silicon edge devices. Through systematic experiments on Qwen3.5 series (0.8B, 2B, 4B parameters) and Phi-3 Mini (3.82B parameters) using Apple Silicon M4 Pro with 24GB unified memory, we demonstrate:

\begin{enumerate}[leftmargin=*,nosep]
    \item \textbf{Framework comparison}: MLX INT4 quantization provides dramatic speedups across all models, with Qwen3.5-0.8B achieving 269.88 tokens/second (9.9$\times$ faster than Ollama), Qwen3.5-2B at 180.48 tokens/second (10.0$\times$ faster), Qwen3.5-4B at 80.69 tokens/second (15.2$\times$ faster), and Phi-3 Mini at 96.53 tokens/second (7.4$\times$ faster than Ollama).
    
    \item \textbf{Memory efficiency}: All models fit comfortably within 8GB memory with MLX INT4 quantization, with memory footprints ranging from 0.48GB (Qwen3.5-0.8B) to 2.47GB (Qwen3.5-4B), enabling deployment on consumer devices.
    
    \item \textbf{Practical deployment}: MLX framework selection is critical for Apple Silicon, providing 7.4$\times$ to 15.2$\times$ speedup over Ollama depending on model size, with larger models showing greater benefits from hardware optimization.
\end{enumerate}
As data center energy consumption approaches 945 TWh annually, edge deployment offers a sustainable alternative. Our experiments demonstrate that framework selection dramatically impacts performance on Apple Silicon, with MLX providing up to 9.9$\times$ speedup over Ollama. For resource-constrained edge devices requiring real-time inference, Qwen3.5-0.8B with MLX INT4 quantization achieves exceptional throughput at 269.88 tokens/second with only 0.48GB memory.

Future work should include evaluation of Qwen3.5 thinking mode for complex reasoning tasks, cross-architecture comparison on NVIDIA and mobile hardware, and integration of carbon-aware scheduling for truly green inference.

% =============================================================================
% ACKNOWLEDGMENTS
% =============================================================================
\section*{Acknowledgments}

We thank the anonymous reviewers for their valuable feedback. This research was supported by Chitkara University's Research Infrastructure. We acknowledge the developers of Hugging Face Transformers, CodeCarbon, and the Apple MLX team for making their tools openly available.

% =============================================================================
% REFERENCES
% =============================================================================
\bibliographystyle{plainnat}
\bibliography{references}

% =============================================================================
% APPENDIX
% =============================================================================
\appendix

\section{Detailed Experimental Setup}

\subsection{Hardware Configuration}

All experiments were conducted on the following hardware:

\begin{table}[h]
\centering
\caption{Hardware specifications}
\label{tab:hardware}
\begin{tabular}{ll}
\toprule
\textbf{Component} & \textbf{Specification} \\
\midrule
Processor & Apple M4 Pro (14-core CPU, 20-core GPU) \\
Memory & 24GB unified memory \\
Storage & SSD \\
OS & macOS \\
PyTorch & 2.7.1 with MPS backend \\
\bottomrule
\end{tabular}
\end{table}



\section{Additional Results}



\subsection{Generation Quality Examples}

Figure \ref{fig:examples} shows generation examples from Qwen3.5-0.8B via Ollama.

\begin{figure}[h]
\centering
\fbox{
\begin{minipage}{0.95\linewidth}
\footnotesize
\textbf{Prompt:} Explain the concept of machine learning in simple terms.

\vspace{0.2cm}
\textbf{Qwen3.5-0.8B:} Machine learning is like teaching a computer to learn from examples, just like how you might teach a child to recognize cats by showing them many pictures of cats. The computer looks at data, finds patterns, and then uses those patterns to make predictions about new information it hasn't seen before.
\end{minipage}
}
\caption{Generation quality example from Qwen3.5-0.8B on Apple Silicon M4 Pro via Ollama with Q4\_K\_M quantization.}
\label{fig:examples}
\end{figure}

\section{Reproducibility}

All code, model checkpoints, and experimental logs are available at: \url{https://github.com/harshpreet931/researchPaper-quantization}

The repository includes:
\begin{itemize}[leftmargin=*,nosep]
    \item Quantization scripts for all methods
    \item Evaluation pipeline for efficiency benchmarks
    \item Energy measurement integration with CodeCarbon
    \item Jupyter notebooks reproducing all figures and tables
    \item Docker container for reproducible environment
\end{itemize}

\end{document}
